\chapter{NP-Completeness, Reductions \& Approximation (Module 5)}
\label{ch:np_completeness}

\section{Computational Complexity Foundations}
Computational complexity categorizes problems by the asymptotic resources required for their solution~\cite{clrs2009,karp1972}:
\begin{itemize}
    \item \textbf{Class P:} Decision problems decidable on a deterministic Turing machine in polynomial time $O(n^c)$ for constant $c$.
    \item \textbf{Class NP:} Decision problems verifiable in polynomial time given a polynomial-length certificate (witness).
    \item \textbf{Polynomial-Time Reduction ($\le_p$):} A problem $A$ reduces to $B$ ($A \le_p B$) if an efficient subroutine for $B$ solves $A$ with polynomial-time overhead. If $B \in P$, then $A \in P$.
    \item \textbf{NP-Completeness:} A language $L$ is NP-complete if $L \in \text{NP}$ and $\forall L' \in \text{NP}, \ L' \le_p L$. The Cook-Levin Theorem (1971)~\cite{cook1971} established that the Boolean Satisfiability problem (SAT) is NP-complete.
\end{itemize}

% ------------------------------------------------------------------------------
\section{DPLL Boolean Satisfiability (SAT) Solver}
\subsection{Algorithm Architecture \& Invariants}
A Boolean formula in Conjunctive Normal Form (CNF) is a conjunction of clauses: $\Phi = C_1 \land C_2 \land \dots \land C_m$, where each clause $C_i = (l_{i, 1} \lor \dots \lor l_{i, k})$ is a disjunction of literals.
The Davis-Putnam-Logemann-Loveland (DPLL) algorithm decides satisfiability via depth-first search augmented by two deterministic simplification rules:
\begin{enumerate}
    \item \textbf{Unit Propagation:} If a clause contains exactly one unassigned literal $l$, the formula can be satisfied only by assigning $l \gets \textbf{true}$. Propagating this forced assignment simplifies all dependent clauses.
    \item \textbf{Pure Literal Elimination:} If a variable appears exclusively with positive (or negative) polarity across all unsatisfied clauses, assign it to satisfy all containing clauses.
    \item \textbf{Splitting / Backtracking:} If no unit clauses remain, select an unassigned variable $x$, branch on $x \gets \textbf{true}$, and backtrack to $x \gets \textbf{false}$ if a contradiction (an empty clause) is derived.
\end{enumerate}

\subsection{Short Illustrative Example (Independent from Project)}
Consider $\Phi = (x_1 \lor x_2) \land (\neg x_1 \lor x_2) \land (x_1 \lor \neg x_2)$.
\begin{enumerate}
    \item Branch on $x_1 \gets \textbf{true}$:
    \begin{itemize}
        \item Clause $(x_1 \lor x_2)$ is satisfied.
        \item Clause $(\neg x_1 \lor x_2)$ reduces to unit clause $(x_2)$.
        \item By Unit Propagation: $x_2 \gets \textbf{true}$.
        \item Clause $(x_1 \lor \neg x_2)$ becomes $(\textbf{true} \lor \textbf{false}) = \textbf{true}$.
    \end{itemize}
    \item All clauses satisfied. Formula is \textbf{SATISFIABLE} under $\{x_1 = \textbf{true}, x_2 = \textbf{true}\}$.
\end{enumerate}

\subsection{EduTrack Domain Application: Examination Timetabling Constraints}
In EduTrack (\texttt{DpllSatSolver.java}, \texttt{ExamSchedulingService.java}), the DPLL engine solves examination timetabling. Let variable $v(c, s)$ denote course $c$ assigned to exam slot $s$:
\begin{align}
\text{Coverage:} &\quad \forall c, \ \bigvee_{s} v(c, s) \nonumber \\
\text{Singular Slot:} &\quad \forall c, s_1 < s_2, \ (\neg v(c, s_1) \lor \neg v(c, s_2)) \nonumber \\
\text{Student Conflict:} &\quad \forall \text{conflicting } (c_1, c_2), \forall s, \ (\neg v(c_1, s) \lor \neg v(c_2, s))
\end{align}
The DPLL solver verifies whether a conflict-free examination schedule exists within available slots.

% ------------------------------------------------------------------------------
\section{The Karp Reduction Chain: 3-SAT to Vertex Cover}
To demonstrate computational intractability from first principles, EduTrack implements the complete polynomial reduction pipeline:
\begin{equation}
\text{3-SAT} \ \le_p \ \text{CLIQUE} \ \le_p \ \text{INDEPENDENT-SET} \ \le_p \ \text{VERTEX-COVER}
\end{equation}

\subsection{Step 1: 3-SAT to CLIQUE Reduction}
Given a 3-CNF formula $\Phi = C_1 \land \dots \land C_m$ where $C_r = (l_{r, 1} \lor l_{r, 2} \lor l_{r, 3})$:
\begin{enumerate}
    \item Construct a graph $G = (V, E)$ with $|V| = 3m$ vertices, one for each literal in each clause: $v_{r, j}$ represents literal $l_{r, j}$.
    \item Add an undirected edge $(v_{r, j}, v_{s, k})$ if and only if:
    \begin{equation}
    r \ne s \quad \text{(different clauses)} \quad \textbf{and} \quad l_{r, j} \ne \neg l_{s, k} \quad \text{(non-contradictory literals)}
    \end{equation}
\end{enumerate}
\begin{theorem}
$\Phi$ is satisfiable if and only if $G$ contains a clique of size $m$.
\end{theorem}
\begin{proof}
If $\Phi$ is satisfied, choose one true literal per clause. The corresponding $m$ vertices belong to different clauses and cannot contradict, forming an $m$-clique. Conversely, an $m$-clique picks exactly one vertex per clause with no contradictory assignments.
\end{proof}

\subsection{Step 2: CLIQUE to INDEPENDENT SET Reduction}
Given graph $G = (V, E)$ and target clique size $k$, construct the complement graph $\overline{G} = (V, \overline{E})$ where $(u, v) \in \overline{E} \iff (u, v) \notin E$.
A subset of vertices $S \subseteq V$ forms a clique in $G$ if and only if every pair in $S$ shares an edge in $G$, meaning no pair in $S$ shares an edge in $\overline{G}$. Hence, $S$ is an \textbf{Independent Set of size $k$ in $\overline{G}$}.

\subsection{Step 3: INDEPENDENT SET to VERTEX COVER Reduction}
\begin{theorem}[Complementary Duality Theorem~\cite{erickson2019}]
In any graph $H = (V, E)$, a subset $S \subseteq V$ is an Independent Set if and only if its complement $V \setminus S$ is a Vertex Cover.
\end{theorem}
\begin{proof}
Let $S$ be an independent set. For any edge $(u, v) \in E$, both $u$ and $v$ cannot belong to $S$. Therefore, at least one of $\{u, v\}$ belongs to $V \setminus S$, making $V \setminus S$ a vertex cover. The converse follows symmetrically.
\end{proof}
Therefore, $\overline{G}$ has an Independent Set of size $k \iff \overline{G}$ has a Vertex Cover of size $|V| - k$.

\subsection{EduTrack Experimental Demonstration}
In EduTrack (\texttt{SatToCliqueReduction.java}, \texttt{CliqueToISReduction.java}, \texttt{ISToVertexCover.java}), a 3-clause timetable constraint formula:
\begin{equation}
\Phi = (x_1 \lor x_2 \lor \neg x_3) \land (\neg x_1 \lor x_3 \lor x_4) \land (\neg x_2 \lor \neg x_3 \lor \neg x_4)
\end{equation}
is transformed across the reduction pipeline:
\begin{itemize}
    \item 3-SAT $\to$ CLIQUE: Produces $G$ with $|V| = 9$, $|E| = 22$, target clique $k = 3$.
    \item CLIQUE $\to$ IS: Produces complement $\overline{G}$ with $|V| = 9$, $|\overline{E}| = 14$, target IS $k = 3$.
    \item IS $\to$ VC: Produces dual Vertex Cover of target size $|V| - k = 9 - 3 = \mathbf{6}$.
\end{itemize}

% ------------------------------------------------------------------------------
\section{Vertex Cover 2-Approximation via Maximal Matching}
\subsection{Algorithm Formulation \& Provable Approximation Guarantee}
Because Minimum Vertex Cover is NP-hard, large-scale instances require polynomial-time approximation algorithms. EduTrack implements the classical maximal matching 2-approximation:
\begin{algorithm}[H]
\caption{Vertex Cover 2-Approximation via Maximal Matching}
\begin{algorithmic}[1]
\Function{ApproxVertexCover}{$G = (V, E)$}
    \State $C \gets \emptyset$ \Comment{Vertex Cover Set}
    \State $E' \gets E$
    \While{$E' \ne \emptyset$}
        \State Pick an arbitrary edge $(u, v) \in E'$
        \State $C \gets C \cup \{u, v\}$
        \State Remove from $E'$ all edges incident to $u$ or $v$
    \EndWhile
    \State \Return $C$
\EndFunction
\end{algorithmic}
\end{algorithm}

\begin{theorem}
Algorithm~\ref{alg:approx_vc} is a polynomial-time 2-approximation algorithm for Minimum Vertex Cover:
\begin{equation}
|C| \le 2 \cdot |\text{OPT}|
\end{equation}
\end{theorem}
\begin{proof}
Let $M$ denote the set of edges selected in line 5. By construction, no two edges in $M$ share an endpoint; hence $M$ is a valid matching, and $|C| = 2|M|$. To cover the edges in $M$, any valid vertex cover (including the optimal cover $\text{OPT}$) must pick at least one distinct vertex per edge in $M$. Thus, $|\text{OPT}| \ge |M|$. Combining these bounds:
\begin{equation}
|C| = 2|M| \le 2 \cdot |\text{OPT}|
\end{equation}
\end{proof}

\subsection{EduTrack Domain Application: Exam Proctor Minimization}
In EduTrack (\texttt{VertexCover2Approx.java}), each vertex represents an exam course, and an edge represents a scheduling conflict (shared students). Assigning proctors/invigilators to the vertex cover guarantees that every conflict edge has an assigned proctor to oversee the conflict, with the total number of proctors guaranteed to be within $2\times$ of the theoretical optimum.

\subsection{Screenshot \& Verification Specifications}
\begin{figure}[H]
    \centering
    \fbox{\begin{minipage}{0.85\textwidth}
        \centering
        \vspace{1.2cm}
        \textbf{[PLACEHOLDER: Screenshot 11 - DPLL SAT and 2-Approx Vertex Cover]}\\[0.2cm]
        \textit{Capture the GUI NP-Completeness tab displaying:}\\
        \textit{1. DPLL SAT Solver truth assignments for exam timetable constraints.}\\
        \textit{2. Karp Reduction step-by-step translation from 3-SAT to Vertex Cover.}\\
        \textit{3. 2-Approximation Vertex Cover proctor allocation and maximal matching size.}
        \vspace{1.2cm}
    \end{minipage}}
    \caption{NP-Completeness Reductions and Approximation in EduTrack UI}
    \label{fig:screenshot_np}
\end{figure}
